% =========================================================================
\section{Introducción}\label{sec:introduccion}
% =========================================================================
En este PreLab se obtienen, con el modelo ideal del amplificador operacional (AO), las ganancias de
tres circuitos realimentados: el amplificador inversor ($A_V=-10$), el no inversor ($A_V=+11$) y el
inversor generalizado, que con un condensador en la realimentación se comporta como un filtro paso bajo.

Un AO amplifica la diferencia entre sus entradas no inversora ($+$) e inversora ($-$),
$\vout=A\,(\vp-\vm)$, con una ganancia en lazo abierto $A$ muy grande; tiene una resistencia de entrada
$\Rin$ entre ambas entradas y una resistencia de salida $\Rout$. El \emph{AO ideal} cumple
\[
  \Rin\to\infty,\qquad \Rout\to0,\qquad A\to\infty .
\]
Si trabaja en zona lineal gracias a una \emph{realimentación negativa} (una red que lleva parte de la
salida a la entrada inversora), se derivan las dos reglas que usaremos en todos los cálculos:
\begin{itemize}[leftmargin=3em]
  \item[(R1)] \emph{corrientes de entrada nulas}, $i_+=i_-=0$, porque entre las entradas solo hay $\Rin\to\infty$;
  \item[(R2)] \emph{cortocircuito virtual}, $\vp=\vm$, porque la salida es finita ($\abs{\vout}\le\Vsat$) y
        $\vp-\vm=\vout/A\to0$ cuando $A\to\infty$.
\end{itemize}
Además, con $\Rout=0$ la tensión de salida no depende de la carga. Cada circuito se resuelve imponiendo
(R1) y (R2) y aplicando la ley de Kirchhoff de corrientes en el nudo inversor y la ley de Ohm.
